\chapter{Ulnar nerve release}
\index{Ulnar nerve release}

\section*{Definition and Overview}
The ulnar nerve is one of the three primary peripheral nerves supplying the upper extremity, playing a critical role in both the sensory and motor function of the arm and hand. Originating from the medial cord of the brachial plexus and carrying fibers from the C8 and T1 spinal nerve roots, the ulnar nerve travels down the medial aspect of the arm. It passes behind the medial epicondyle of the humerus, through a narrow fibro-osseous tunnel known as the cubital tunnel at the elbow, continues down the medial forearm, and enters the hand through Guyon's canal at the wrist.

Ulnar nerve release refers to the surgical decompression of this nerve along its path, most frequently at the elbow, to treat cubital tunnel syndrome. Cubital tunnel syndrome is the second most common compression neuropathy of the upper limb, surpassed only by carpal tunnel syndrome of the median nerve. Compression, traction, or friction on the ulnar nerve can lead to progressive ischemia, demyelination, and eventual axonal loss. This pathogenetic cascade manifests as numbness, tingling, and paresthesia in the small finger and the ulnar half of the ring finger, accompanied by hand weakness and intrinsic muscle wasting in advanced stages.

The primary objective of ulnar nerve release is to alleviate mechanical pressure on the nerve, thereby restoring microvascular perfusion and facilitating axonal regeneration or remyelination. Surgical intervention is indicated when conservative measures fail to provide relief, or in cases presenting with moderate-to-severe neurological deficits, such as muscle weakness or motor wasting. Understanding the detailed anatomy of the cubital tunnel, including the role of Osborne's ligament (the roof of the tunnel) and the deep aponeurosis of the flexor carpi ulnaris muscle, is paramount for successful decompression. This chapter provides a comprehensive review of the clinical presentation, diagnosis, surgical options, complications, and post-operative management associated with ulnar nerve release.

\section*{Epidemiology}
Compression neuropathies of the ulnar nerve represent a significant clinical and socioeconomic burden. While decompression can be performed at the wrist (Guyon's canal), the vast majority of cases occur at the elbow within the cubital tunnel. The estimated incidence of cubital tunnel syndrome is approximately 25 to 30 cases per 100,000 person-years, making it a highly prevalent upper extremity disorder.

Epidemiological studies indicate that the condition is significantly more common in males than in females, with some literature reporting a male-to-female ratio of nearly 2:1. This gender disparity is hypothesized to be due to differences in subcutaneous fat distribution around the elbow, muscular anatomy, and occupational exposure. The peak age of presentation occurs between the fifth and sixth decades of life, although it can affect individuals of any age, including adolescents and young athletes.

Certain populations and occupations exhibit a heightened risk of developing ulnar nerve compression. Individuals whose occupations require repetitive elbow flexion and extension, prolonged leaning on hard surfaces, or exposure to chronic vibration are particularly vulnerable. Examples include assembly line workers, truck drivers, office workers, computer programmers, and musicians. Furthermore, overhead throwing athletes, such as baseball pitchers, are prone to ulnar nerve pathology due to the extreme valgus stress placed on the elbow joint during the throwing motion.

Co-existing systemic diseases also influence the epidemiology of ulnar neuropathy. Patients with diabetes mellitus, thyroid disorders, chronic renal failure, and obesity exhibit higher rates of compression neuropathy. In diabetics, the nerve's heightened susceptibility to compression is often attributed to subclinical peripheral neuropathy and compromised microvascular circulation, which lowers the threshold for symptomatic nerve compression.

\section*{Aetiology and Pathophysiology}
The ulnar nerve is uniquely vulnerable to compression and traction because of its anatomical course, particularly at the elbow where it passes posterior to the flexion axis. The etiology of ulnar nerve compression is multi-factorial, involving anatomical, mechanical, and systemic factors. The pathophysiological mechanisms are primarily driven by mechanical compression, traction, and friction, which lead to local ischemia, intraneural edema, and focal demyelination.

Under normal conditions, elbow flexion decreases the volume of the cubital tunnel by flattening Osborne's ligament and tightening the deep fascia of the flexor carpi ulnaris. This movement can increase the pressure within the tunnel up to twenty-fold, compressing the nerve. Concurrently, because the ulnar nerve lies posterior to the elbow's axis of rotation, elbow flexion stretches the nerve, increasing longitudinal tension and reducing the cross-sectional area of the nerve itself. Chronic or repetitive exposure to these forces compromises the microvascular blood supply (vasa nervorum) to the nerve, leading to endoneurial hypoxia and edema. If the pressure persists, the nerve undergoes demyelination, followed by axonal degeneration in severe cases.

The key etiological factors and mechanisms include:
\begin{itemize}
  \item \textbf{Mechanical Compression}: The ulnar nerve can be compressed at several distinct anatomical sites around the elbow. The most common site is Osborne's ligament, which spans from the medial epicondyle to the olecranon. Other sites include the arcade of Struthers (a thin fascial band in the distal medial arm), the medial intermuscular septum, the deep aponeurosis of the flexor carpi ulnaris (FCU) muscle, and the exit of the nerve from the FCU muscle belly.
  \item \textbf{Traction Forces}: Repetitive or sustained elbow flexion (greater than 90 degrees) stretches the ulnar nerve. Over time, this traction leads to chronic micro-trauma, intraneural scarring, and loss of nerve elasticity.
  \item \textbf{Frictional Neuritis and Subluxation}: In approximately 15\% to 20\% of the population, the ulnar nerve is hyper-mobile and subluxates or dislocates anteriorly over the medial epicondyle during elbow flexion. This movement causes the nerve to rub against the bony prominence, resulting in chronic friction, mechanical irritation, and inflammatory changes.
  \item \textbf{Trauma and Deformities}: Acute trauma, such as fractures or dislocations of the elbow, can directly injure the ulnar nerve. Additionally, chronic deformities such as cubitus valgus (often a late complication of childhood supracondylar fractures, historically referred to as "tardy ulnar palsy") increase the pathway distance the nerve must travel, putting it under chronic traction.
  \item \textbf{Space-Occupying Lesions}: Local pathology within or adjacent to the cubital tunnel can compress the nerve. Examples include ganglion cysts, lipomas, rheumatoid pannus, osteophytes associated with osteoarthritis, and heterotopic ossification.
  \item \textbf{Systemic and Metabolic Vulnerabilities}: Systemic conditions such as diabetes mellitus, hypothyroidism, chronic alcoholism, and renal failure alter nerve metabolism and microvascular flow, rendering the nerve more susceptible to compression at lower pressure thresholds.
\end{itemize}

\section*{Clinical Features}
The clinical presentation of ulnar nerve compression varies depending on the severity and duration of the condition. Symptoms generally progress from intermittent sensory disturbances to constant numbness, followed by motor weakness and muscle atrophy in the hand.

\begin{itemize}
  \item \textbf{Sensory Disturbances}: The earliest and most common symptom is paresthesia, numbness, and tingling in the ulnar nerve distribution, which includes the little finger and the medial (ulnar) half of the ring finger. These symptoms are frequently worse at night, often because patients sleep with their elbows flexed. Activities that require prolonged elbow flexion, such as holding a telephone to the ear, driving, or typing, also exacerbate the symptoms.
  \item \textbf{Pain}: Patients may experience a dull, aching pain or discomfort over the medial aspect of the elbow, which can radiate distally down the forearm or proximally into the arm. Direct pressure on the medial epicondyle (the "funny bone") can cause sharp, radiating pain.
  \item \textbf{Motor Weakness and Clumsiness}: As the compression worsens, patients notice a loss of manual dexterity. They may drop objects, have difficulty opening jars, or struggle with fine motor tasks such as buttoning shirts, tying shoelaces, or playing musical instruments. This weakness is due to the loss of innervation to the intrinsic muscles of the hand.
  \item \textbf{Muscle Atrophy and Deformity}: In severe, chronic cases, visible wasting of the hand muscles becomes apparent. Atrophy is most prominent in the first dorsal interosseous muscle (visible as a depression in the web space between the thumb and index finger), the hypothenar eminence, and the interossei muscles. In advanced stages, the patient may develop an "ulnar claw hand" (Duchenne's sign), characterized by hyperextension of the metacarpophalangeal joints and flexion of the interphalangeal joints of the ring and little fingers, resulting from the loss of function of the third and fourth lumbrical muscles.
  \item \textbf{Provocative Clinical Signs}: During physical examination, several clinical tests can reproduce or confirm ulnar nerve dysfunction. Tinel's sign is positive when light percussion over the ulnar nerve at the cubital tunnel elicits paresthesia in the ulnar digits. The Elbow Flexion Test involves holding the patient's elbow in full flexion with the wrist in neutral for up to sixty seconds; a positive test reproduces numbness or tingling. Wartenberg's sign is present when the patient is unable to adduct the little finger against the ring finger, caused by weakness of the third palmar interosseous muscle. Froment's sign is demonstrated during a key pinch test: when asked to pinch a piece of paper between the thumb and index finger, the patient compensates for weak adductor pollicis muscle (ulnar nerve) by flexing the interphalangeal joint of the thumb, which is controlled by the flexor pollicis longus (median nerve). Jeanne's sign, which is hyperextension of the thumb metacarpophalangeal joint during pinch, may also be observed.
\end{itemize}

\section*{Differential Diagnosis}
Because symptoms of numbness and weakness in the hand can originate from various locations along the neurological pathway, a detailed differential diagnosis is essential. The clinician must rule out proximal lesions in the cervical spine and brachial plexus, as well as alternative distal neuropathies.

\begin{itemize}
  \item \textbf{C8--T1 Cervical Radiculopathy}: Compression of the C8 or T1 nerve roots in the cervical spine (due to a herniated disc or foraminal stenosis) can mimic the sensory symptoms of ulnar neuropathy. However, cervical radiculopathy typically presents with neck pain, stiffness, and symptoms that radiate from the neck down the arm. Furthermore, sensory deficits in radiculopathy often extend proximal to the wrist, and motor weakness may involve non-ulnar muscles that share C8--T1 innervation, such as the abductor pollicis brevis (median nerve).
  \item \textbf{Thoracic Outlet Syndrome (TOS)}: This condition involves the compression of the lower trunk of the brachial plexus (which contains C8--T1 fibers) as it passes through the thoracic outlet. TOS can cause pain, numbness, and weakness in the ulnar distribution of the hand. It is often differentiated by proximal shoulder or clavicular pain, symptoms exacerbated by overhead activities, and vascular signs such as arm swelling, coolness, or diminished pulses (if the subclavian vessels are also compressed).
  \item \textbf{Pancoast Tumor}: A tumor at the apex of the lung can invade the lower trunk of the brachial plexus, producing severe, unremitting pain along the ulnar side of the arm and hand, along with intrinsic muscle wasting. It is frequently accompanied by Horner's syndrome (miosis, ptosis, anhidrosis) due to involvement of the sympathetic chain.
  \item \textbf{Guyon's Canal Syndrome (Ulnar Tunnel Syndrome)}: This is compression of the ulnar nerve at the wrist rather than the elbow. It can be distinguished because Guyon's canal syndrome spares the sensation on the dorsum of the ulnar digits. The dorsal cutaneous branch of the ulnar nerve branches off approximately 5 to 8 centimeters proximal to the wrist and does not enter Guyon's canal; therefore, dorsal hand sensation remains intact in wrist-level compression but is compromised in elbow-level compression.
  \item \textbf{Medial Epicondylitis (Golfer's Elbow)}: This is an overuse tendinopathy of the common flexor tendon origin at the medial epicondyle. While it causes pain in the same regional area as cubital tunnel syndrome, it is characterized by localized tenderness and pain with resisted wrist flexion, and it does not present with numbness, tingling, or muscle weakness in the hand.
  \item \textbf{Double Crush Syndrome}: This refers to the co-existence of two separate compressions along the same nerve pathway (for example, a mild cervical radiculopathy combined with a mild cubital tunnel syndrome). The proximal compression renders the nerve more vulnerable to distal compression. Recognizing double crush syndrome is vital, as treating only one compression site may lead to incomplete resolution of symptoms.
  \item \textbf{Motor Neuron Disease (Amyotrophic Lateral Sclerosis)}: Early ALS can present with weakness and wasting of the intrinsic hand muscles. However, ALS does not cause sensory symptoms (numbness or pain) and is characterized by progressive, widespread muscle wasting, fasciculations, and upper motor neuron signs such as hyperreflexia and spasticity.
\end{itemize}

\section*{When to Seek Medical Care}
Early diagnosis and intervention are critical to preventing irreversible nerve damage. Patients should monitor their symptoms and seek professional medical evaluation if they experience persistent or worsening signs of ulnar nerve compression.

A medical consultation is recommended if:
\begin{itemize}
  \item Numbness, tingling, or pain in the hand and fingers becomes frequent, constant, or interferes with sleep.
  \item There is a noticeable decrease in hand strength, such as difficulty opening bottles, holding objects, or performing fine motor tasks.
  \item Visible muscle wasting (atrophy) is observed in the hand, particularly in the web space between the thumb and index finger.
  \item Conservative measures (such as activity modification or night splinting) fail to improve symptoms after several weeks.
\end{itemize}

Immediate or emergency medical attention should be sought if:
\begin{itemize}
  \item The loss of sensation or motor function occurs suddenly, particularly following an acute injury, fracture, or dislocation of the elbow.
  \item Symptoms of numbness or weakness in the arm or hand are accompanied by sudden difficulty speaking, facial drooping, confusion, dizziness, or weakness on one side of the body, which are warning signs of a stroke.
  \item The arm pain is accompanied by chest tightness, shortness of breath, sweating, or pain radiating to the jaw or neck, which may indicate a myocardial infarction (heart attack).
\end{itemize}
For life-threatening emergencies, such as suspected stroke or heart attack, do not delay. Call emergency services immediately. In many countries, call local emergency services; in the United States, call 911; in the United Kingdom, call 999.

\section*{Diagnosis and Investigations}
The diagnosis of ulnar nerve compression is primarily clinical, based on a detailed history and physical examination, and is confirmed using electrodiagnostic and imaging studies.

\begin{itemize}
  \item \textbf{Clinical Examination}: A thorough neurological examination is the cornerstone of diagnosis. Sensory testing includes assessing light touch, pain sensation, and two-point discrimination in the ulnar digits compared to the radial and median nerve distributions. Motor testing evaluates the strength of ulnar-innervated muscles, including the flexor carpi ulnaris (FCU), the flexor digitorum profundus (FDP) of the ring and little fingers, and the intrinsic hand muscles (abductor digiti minimi, interossei, and adductor pollicis). Provocative maneuvers, including Tinel's sign at the elbow, the Elbow Flexion Test, and checking for Froment's, Wartenberg's, and Jeanne's signs, are performed to localize the site of compression.
  \item \textbf{Electrodiagnostic Studies (EMG and NCS)}: Electromyography (EMG) and Nerve Conduction Studies (NCS) are essential to confirm the diagnosis, grade the severity of the neuropathy, and rule out other sites of compression. Key electromyographic findings in cubital tunnel syndrome include a reduction in ulnar nerve conduction velocity across the elbow segment to less than 50 meters per second. Conduction block or temporal dispersion across the elbow further localizes the pathology. EMG of the ulnar-innervated muscles (such as the first dorsal interosseous) may show signs of active denervation, such as fibrillation potentials and positive sharp waves, or signs of chronic reinnervation, such as polyphasic motor unit action potentials.
  \item \textbf{High-Resolution Musculoskeletal Ultrasonography}: Ultrasound has emerged as a valuable diagnostic tool, providing detailed anatomical information. In cubital tunnel syndrome, ultrasound can demonstrate swelling of the ulnar nerve proximal to the site of compression. A cross-sectional area of the ulnar nerve at the medial epicondyle exceeding 9 to 10 square millimeters is highly suggestive of pathology. Ultrasound is also useful for dynamic evaluation, allowing the clinician to visualize ulnar nerve subluxation or dislocation over the medial epicondyle during elbow flexion, as well as detecting space-occupying lesions such as ganglion cysts or accessory muscles.
  \item \textbf{Magnetic Resonance Imaging (MRI)}: While not routinely required for all patients, MRI of the cervical spine or elbow is indicated in atypical presentations or when proximal compression (radiculopathy, thoracic outlet syndrome) is suspected. Elbow MRI can reveal localized ulnar nerve swelling, increased signal intensity on T2-weighted sequences (indicating intraneural edema), and denervation changes (such as muscle edema or atrophy) in the ulnar-innervated forearm and hand muscles. It is also excellent for identifying deep-seated space-occupying lesions.
  \item \textbf{Plain Radiographs (X-rays)}: Standard anteroposterior and lateral radiographs of the elbow are performed to evaluate for bony abnormalities that could compress the nerve. These include osteophytes from osteoarthritis, heterotopic ossification, calcified ligaments, cubitus valgus deformity, or malunion from old fractures.
\end{itemize}

\section*{Management}
Management of ulnar nerve compression is tailored to the severity of symptoms, the presence of objective neurological deficits, and the patient's lifestyle and occupational needs. Treatment ranges from conservative, non-operative therapies for mild cases to surgical intervention for moderate-to-severe or progressive neuropathy.

Conservative Management:
\begin{itemize}
  \item \textbf{Activity Modification}: Patients are advised to avoid activities that require repetitive or prolonged elbow flexion (greater than 90 degrees) and to avoid leaning directly on their elbows. Ergonomic adjustments at work, such as using a headset for phone calls, modifying keyboard height, and adjusting chair armrests, can significantly reduce nerve compression.
  \item \textbf{Night Splinting}: An elbow extension splint is prescribed for night use, holding the elbow in a comfortable position of approximately 30 to 45 degrees of flexion. This prevents prolonged, deep flexion during sleep, reducing tension and pressure on the nerve.
  \item \textbf{Nerve Gliding Exercises}: Physical and occupational therapists teach nerve gliding (or sliding) exercises. These movements encourage the ulnar nerve to glide smoothly through the cubital tunnel and Guyon's canal, reducing adhesions and tension within the nerve bed.
  \item \textbf{Pharmacotherapy}: Non-steroidal anti-inflammatory drugs (NSAIDs), such as ibuprofen or naproxen, may be used for short-term pain management to reduce localized inflammation. Oral corticosteroids or corticosteroid injections are generally not recommended in the cubital tunnel due to limited efficacy and the risk of direct needle injury to the nerve.
\end{itemize}

Surgical Management (Ulnar Nerve Release):
\begin{itemize}
  \item \textbf{In Situ Cubital Tunnel Release}: This is the most common and straightforward surgical procedure. Through an incision over the medial aspect of the elbow, the surgeon decompresses the nerve by dividing Osborne's ligament and the deep aponeurosis of the flexor carpi ulnaris muscle. The nerve is left in its native anatomical position behind the medial epicondyle. This technique has a lower rate of complications and a faster recovery time, but it may not be suitable if the nerve is prone to subluxating post-release.
  \item \textbf{Anterior Transposition}: If the ulnar nerve is unstable (subluxates over the medial epicondyle during flexion) or if there is severe deformity or scarring, the surgeon may decompress and move (transpose) the nerve anterior to the medial epicondyle. The nerve can be placed in one of three positions: subcutaneous (under the skin and fat), intramuscular (within the pronator-flexor muscle mass), or submuscular (beneath the flexor-pronator muscle origin, which is the most complex but provides the most protected environment). Anterior transposition eliminates traction on the nerve during elbow flexion.
  \item \textbf{Medial Epicondylectomy}: In this procedure, the nerve is decompressed in situ, and the medial epicondyle bone prominence is surgically removed. This removes the pulley around which the ulnar nerve must bend, reducing traction and compression.
  \item \textbf{Endoscopic Ulnar Nerve Release}: A minimally invasive technique that uses a smaller incision and an endoscope to visualize and divide the compressing structures. This method aims to minimize soft tissue dissection, post-operative pain, and scarring, allowing for quicker rehabilitation.
\end{itemize}

\section*{Complications}
While ulnar nerve release is generally a safe and highly successful procedure, as with any surgery, there are inherent risks and potential complications.

\begin{itemize}
  \item \textbf{Injury to the Medial Antebrachial Cutaneous Nerve (MCN)}: The branches of the MCN cross the medial aspect of the elbow near the surgical site. Damage to these branches can cause numbness or chronic, burning pain on the medial side of the forearm, and can lead to the formation of a painful neuroma, which is difficult to treat.
  \item \textbf{Incomplete Decompression}: If the surgeon fails to release all potential sites of compression (such as the arcade of Struthers, the medial intermuscular septum, or the distal FCU fascia), the patient may experience persistent symptoms after surgery.
  \item \textbf{Perineural Scarring}: Post-operative scarring around the decompressed nerve can cause the nerve to become tethered to the surrounding tissues, leading to recurrent pain and paresthesia.
  \item \textbf{Ulnar Nerve Subluxation}: Following an in situ release, the ulnar nerve may become hyper-mobile and begin to slide over the medial epicondyle during elbow flexion. This subluxation can cause mechanical irritation and recurrent symptoms, occasionally requiring a revision surgery to perform an anterior transposition.
  \item \textbf{Infection and Wound Complications}: Surgical site infections, hematomas, or delayed wound healing can occur along the medial elbow incision, particularly in patients with risk factors like diabetes or smoking.
  \item \textbf{Elbow Stiffness}: Prolonged immobilization after surgery can lead to joint stiffness and a temporary loss of full extension or flexion. Early, controlled range-of-motion exercises are crucial to prevent this.
  \item \textbf{Complex Regional Pain Syndrome (CRPS)}: A rare but severe complication characterized by chronic, disproportionate pain, swelling, temperature changes, and sweating in the affected limb, requiring intensive pain management and physical therapy.
\end{itemize}

\section*{Prognosis and Outcomes}
The prognosis following ulnar nerve release is generally favorable, with the majority of patients achieving significant symptom relief and functional improvement. The success of the procedure depends heavily on the preoperative severity of the nerve compression, the duration of symptoms, and the patient's general health.

For patients with mild to moderate cubital tunnel syndrome (who experience intermittent numbness and no muscle wasting), conservative treatment resolves symptoms in approximately 50\% to 60\% of cases. For those who undergo surgical decompression, success rates range from 85\% to 90\%. These patients typically experience rapid relief from nighttime paresthesia and pain, followed by a gradual improvement in daytime numbness.

In contrast, the prognosis is more guarded for patients presenting with severe, chronic compression characterized by constant numbness, muscle atrophy, and clawing. While surgical release is successful in halting the progression of the neuropathy, complete recovery of muscle bulk and fine motor strength is less common. Nerves regenerate at a rate of approximately 1 millimeter per day, meaning sensory and motor recovery is a slow process that can take up to twelve to eighteen months post-surgery.

Factors that negatively impact the prognosis include advanced age, long duration of symptoms (greater than one year), severe preoperative muscle wasting, poorly controlled diabetes mellitus, and smoking, which compromises microvascular healing. Adherence to a structured post-operative physical or occupational therapy program is a critical determinant of the final functional outcome.

\section*{Prevention and Self-Care}
Preventing ulnar nerve compression and managing mild, early-stage symptoms involve minimizing mechanical stress, friction, and traction on the nerve through ergonomic changes and lifestyle adjustments.

\begin{itemize}
  \item \textbf{Ergonomic Adjustments}: When working at a desk or computer, adjust the height of the chair and keyboard so that the elbows are bent at an angle of 90 degrees or less. Avoid resting the elbows on hard desk surfaces or chair armrests; use padded armrest covers if necessary. Use a hands-free headset or speakerphone for phone calls to avoid prolonged elbow flexion.
  \item \textbf{Regular Breaks and Stretching}: Incorporate frequent breaks into tasks that require repetitive arm movements. Straighten the elbows and perform gentle stretching exercises for the wrists and arms to relieve tension on the nerve.
  \item \textbf{Night Protection}: To prevent deep elbow flexion during sleep, wear a commercial elbow splint or wrap a rolled towel around the elbow, securing it with tape or straps. This keeps the elbow in a slightly extended position (30 to 45 degrees of flexion) and prevents the patient from sleeping with their elbows fully bent or tucked under their head.
  \item \textbf{Avoid Direct Pressure}: Avoid leaning the elbows on the car door or center console while driving, and refrain from resting the elbows on hard surfaces during daily activities.
  \item \textbf{Maintain Metabolic Health}: Manage chronic conditions such as diabetes and hypothyroidism through proper medical care, diet, and exercise. Controlling blood sugar levels is particularly important, as hyperglycemia increases the susceptibility of peripheral nerves to compression injuries.
  \item \textbf{Proper Athletic Technique}: Throwing athletes and weightlifters should work with coaches or physical therapists to ensure proper technique, reducing valgus stress on the elbow joint and minimizing micro-trauma to the ulnar nerve.
\end{itemize}

\section*{Sources and Further Reading}
For reliable, evidence-based information regarding ulnar nerve compression, cubital tunnel syndrome, and surgical decompression, refer to the following resources:

\begin{itemize}
  \item American Academy of Orthopaedic Surgeons (AAOS) -- OrthoInfo:
  https://orthoinfo.aaos.org
  \item American Society for Surgery of the Hand (ASSH) -- HandCare:
  https://www.assh.org
  \item Mayo Clinic -- Cubital Tunnel Syndrome:
  https://www.mayoclinic.org
  \item Cleveland Clinic -- Ulnar Nerve Entrapment:
  https://my.clevelandclinic.org
  \item National Institute of Neurological Disorders and Stroke (NINDS):
  https://www.ninds.nih.gov
  \item Johns Hopkins Medicine -- Cubital Tunnel Syndrome:
  https://www.hopkinsmedicine.org
\end{itemize}